\section{Maraschino Cherries: Distilled Pairwise Alignments}
\label{sec:maraschino}

``Cherries'' are pairwise training examples chosen, in place of full phylogenetically-annotated multiple sequence alignments,
as a composite likelihood approximation to the full phylogenetic likelihood \cite{Prillo2022CherryML}.

``Maraschino Cherries'' are order-1 counts tensors that summarize the adjacency statistics of such pairwise alignments.
MixDom's Maraschino Cherries generalize CherryML to include context-dependent substitution and indel patterns \cite{Prillo2022CherryML}.

The Maraschino pipeline has two phases.
First (Section~\ref{sec:maraschino-counts}), pairwise alignments are reduced to fixed-shape \emph{cherry-count} tensors that aggregate adjacency statistics binned by divergence time.
Second (Section~\ref{sec:maraschino-fit}), the parameters of the MixDom Pair HMM (Section~\ref{sec:nestpairhmm}) are estimated by maximizing the cherry-count log-likelihood: a composite likelihood that scores each adjacency in each time bin under the latent-marginalized collapsed Pair HMM transition matrix $\transnest(\theta,\evoltime)$ derived in Section~\ref{sec:nestpairhmm}.
After fitting, the MixDom model is then distilled to compact order-1 machines (an HMM and a WFST) suitable for use in tree algorithms (Sections~\ref{sec:order1hmm} and \ref{sec:order1trans}).

\subsection{Cherry-count summary statistics}
\label{sec:maraschino-counts}

The input to the Maraschino fitter is a precomputed tensor of pairwise
adjacency counts. For each multiple sequence alignment, sibling pairs are
extracted, gapped columns dropped, and the resulting pairwise alignment
classified into \emph{adjacency contexts}: for every pair of consecutive
non-empty alignment columns, we record the source column type
($\sta$, $\mat$, $\ins$, $\del$), the destination column type
($\mat$, $\ins$, $\del$, $\fin$), and the ancestor/descendant
characters in each.

The pairwise $p$-distance of each cherry is converted to an estimated
divergence time $\evoltime$, and $\evoltime$ is discretized into
$n_\tau$ geometric bins $\{\evoltime_1,\ldots,\evoltime_{n_\tau}\}$ with
representative bin centres $\bar\evoltime_b$.
Counts are accumulated per bin into the following tensors over the
amino-acid alphabet $\alphabet$ ($|\alphabet|=20$), with extended
vocabulary $\alphabet \cup \{\sta, \fin\}$ for boundary positions:

\begin{center}
\begin{tabular}{lll}
\hline
\textbf{Tensor} & \textbf{Shape} & \textbf{Meaning} \\
\hline
$B$ & $n_\tau \times (|\alphabet|+2)^2$
& Singlet bigrams (incl.\ $\sta$/$\fin$) \\
$C^{\mat\mat}$ & $n_\tau \times |\alphabet|^4$
& Match$\to$Match: $(\anctok,\destok,\anctok',\destok')$ \\
$C^{\mat\ins}$ & $n_\tau \times |\alphabet|^3$
& Match$\to$Insert: $(\anctok,\destok,\destok')$ \\
$C^{\mat\del}$ & $n_\tau \times |\alphabet|^3$
& Match$\to$Delete: $(\anctok,\destok,\anctok')$ \\
$C^{\ins\mat}$ & $n_\tau \times |\alphabet|^3$
& Insert$\to$Match: $(\destok,\anctok',\destok')$ \\
$C^{\ins\ins}$ & $n_\tau \times |\alphabet|^2$
& Insert$\to$Insert: $(\destok,\destok')$ \\
$C^{\ins\del}$ & $n_\tau \times |\alphabet|^2$
& Insert$\to$Delete: $(\destok,\anctok')$ \\
$C^{\del\mat}$ & $n_\tau \times |\alphabet|^3$
& Delete$\to$Match: $(\anctok,\anctok',\destok')$ \\
$C^{\del\del}$ & $n_\tau \times |\alphabet|^2$
& Delete$\to$Delete: $(\anctok,\anctok')$ \\
$C^{\del\ins}$ & $n_\tau \times |\alphabet|^2$
& Delete$\to$Insert: $(\anctok,\destok')$ \\
$C^{\sta\mat}$ & $n_\tau \times |\alphabet|^2$
& Start$\to$Match: $(\anctok',\destok')$ \\
$C^{\sta\ins}$ & $n_\tau \times |\alphabet|$
& Start$\to$Insert: $(\destok')$ \\
$C^{\sta\del}$ & $n_\tau \times |\alphabet|$
& Start$\to$Delete: $(\anctok')$ \\
$C^{\mat\fin}$ & $n_\tau \times |\alphabet|^2$
& Match$\to$End: $(\anctok,\destok)$ \\
$C^{\ins\fin}$ & $n_\tau \times |\alphabet|$
& Insert$\to$End: $(\destok)$ \\
$C^{\del\fin}$ & $n_\tau \times |\alphabet|$
& Delete$\to$End: $(\anctok)$ \\
$C^{\sta\fin}$ & $n_\tau$
& Start$\to$End (empty alignment) \\
\hline
\end{tabular}
\end{center}

The post-Insert and post-Delete tensors carry only the adjacent emitted
character ($\destok$ for inserts, $\anctok$ for deletes) as context,
not the previous match's full $(\anctok,\destok)$ context: the model
will marginalise its context-rich frequencies down to this reduced
context when computing the likelihood. The Match-to-Match
tensor $C^{\mat\mat}$ is the largest and dominates the parameter
budget. Boundary tensors record alignments that begin or end
in a particular adjacency type.

\subsection{Cherry-count likelihood for the MixDom Pair HMM}
\label{sec:maraschino-fit}

The cherry-count tensors of Section~\ref{sec:maraschino-counts} are
scored against the collapsed MixDom Pair HMM
$\transnest(\theta,\evoltime)$ defined by
equation~(\ref{eq:mixdom_transitions}) of Section~\ref{sec:nestpairhmm}.
The free parameters are
\[
\theta = \big(\insrate_\main, \delrate_\main,\
\{\insrate_\dom, \delrate_\dom\}_{\dom=1}^{\ndom},\
\{\domdist_\dom\}_{\dom=1}^{\ndom},\
\{\fragdist_{\dom\frag}\}_{\dom,\frag},\
\{\ext^{(\dom)}_{\srcfrag\destfrag}\}_{\dom,\srcfrag,\destfrag},\
\{\classdist_{\dom\frag\class}\}_{\dom,\frag,\class},\
\{\eqm^{(\class)}, \exch^{(\class)}\}_{\class=1}^{\nclasses}\big),
\]
parameterised in unconstrained space (log-rates, log-Dirichlet
weights, log-exchangeabilities) so that gradient methods are
unconstrained. Note in particular:
\begin{itemize}[nosep]
\item $\fragdist_{\dom\frag}$ is the per-domain Dirichlet over fragment
  entry types (initial state of the intra-fragment Markov chain).
\item $\ext^{(\dom)}_{\srcfrag\destfrag}$ is a per-domain
  $\nfrag \times \nfrag$ row-stochastic matrix (with row sums $\leq 1$)
  giving the intra-fragment Markov transition probability from
  fragment-type $\srcfrag$ to fragment-type $\destfrag$ within the
  same fragment; the residual mass
  $\notext^{(\dom)}_\srcfrag = 1 - \sum_\destfrag \ext^{(\dom)}_{\srcfrag\destfrag}$
  is the fragment-termination probability. The TKF92 scalar
  self-extension is the special case $\nfrag=1$, with
  $\ext^{(\dom)}_{11}$ playing the role of the TKF92 extension
  probability (the M-step closed forms for both reduce to a single
  binary count split per row of $\ext^{(\dom)}$).
\item $\classdist_{\dom\frag\class}$ is a per-(domain, fragment-type)
  Dirichlet over $\nclasses$ \emph{static} site classes (drawn
  independently per emitted site, with no chain-time class switching).
\item Each site class $\class$ has its own reversible substitution
  model $\subproc(\exch^{(\class)}, \eqm^{(\class)})$ with rate matrix
  $\revsub^{(\class)} = \exch^{(\class)}\,\diag(\eqm^{(\class)})$.
  All rate variation across sites is captured by these per-class GTR
  matrices; there is no separate Yang-style discretized-gamma
  rate-multiplier mechanism.
\end{itemize}

\paragraph{Pair adjacency frequencies.}
For each $\evoltime$-bin centre $\bar\evoltime_b$, the collapsed
$(5\ndom\nfrag+2)$-state Pair HMM transition matrix
$\transnest^{(b)} \equiv \transnest(\theta,\bar\evoltime_b)$
is constructed via the closed form of
equation~(\ref{eq:mixdom_transitions}). The marginal stationary
distribution $\pi^{\text{stat}}_b$ over emitting states is obtained as the
left null vector of $I - \transnest^{(b)}_{\emit\emit}$, where
$\emit$ denotes the set of $5\ndom\nfrag$ emitting states.

For source machine state $u \in \{\sta, \mat, \ins, \del\}$,
destination machine state $v \in \{\mat, \ins, \del, \fin\}$,
and characters $(\anctok, \destok, \anctok', \destok')$ in their
respective slots, the \emph{model-side adjacency frequency}
\[
F^{uv}_b(\anctok,\destok;\anctok',\destok')
=
\sum_{\substack{s \in \emit^{u} \\ s' \in \emit^{v}}}
\big(\pi^{\text{stat}}_{b,s}\big)
\,\emprob_s(\anctok,\destok)
\,\transnest^{(b)}_{s s'}
\,\emprob_{s'}(\anctok',\destok')
\]
sums over latent (domain, fragment-type) realisations of source
and destination collapsed states $s,s'$, weighting by their
stationary probability and emission probabilities.
The emission probability of a Match state $\mat\mat_{\dom\frag}$
emitting $(\anctok,\destok)$ marginalises the static site-class mixture:
\begin{equation}
\emprob_{\mat\mat_{\dom\frag}}(\anctok,\destok)
= \sum_{\class=1}^{\nclasses}
\classdist_{\dom\frag\class}\,
\eqm^{(\class)}_\anctok
\exp(\revsub^{(\class)} \evoltime)_{\anctok\destok},
\label{eq:maraschino-match-emprob}
\end{equation}
and emission probabilities for $\mat\ins, \ins\ins, \mat\del, \del\del$
analogously marginalise the same class mixture but emit only one of
$(\anctok, \destok)$:
\[
\emprob_{\ins\ins_{\dom\frag}}(\destok)
=\emprob_{\mat\ins_{\dom\frag}}(\destok)
= \sum_\class \classdist_{\dom\frag\class}\,
\eqm^{(\class)}_\destok,
\quad
\emprob_{\del\del_{\dom\frag}}(\anctok)
=\emprob_{\mat\del_{\dom\frag}}(\anctok)
= \sum_\class \classdist_{\dom\frag\class}\,
\eqm^{(\class)}_\anctok.
\]
Boundary frequencies $F^{\sta v}_b$, $F^{u\fin}_b$,
$F^{\sta\fin}_b$ replace the corresponding endpoint factor with the
$\sta$ row or $\fin$ column of $\transnest^{(b)}$.
The reduced-context frequencies needed by post-Insert and post-Delete
counts are obtained by marginalisation:
\[
F^{\ins v}_b(\destok;\cdot) = \sum_{\anctok} F^{\ins v}_b(\anctok,\destok;\cdot),
\qquad
F^{\del v}_b(\anctok;\cdot) = \sum_{\destok} F^{\del v}_b(\anctok,\destok;\cdot).
\]
For each context $(u; \anctok,\destok)$ at bin $b$, the row
normalisation constant is
\[
Z^{u}_b(\anctok,\destok)
= \sum_{v} \sum_{\anctok',\destok'} F^{uv}_b(\anctok,\destok;\anctok',\destok'),
\]
where the inner sum runs over the characters carried by the destination
adjacency type, and the post-Insert/post-Delete normalisations use the
reduced-context frequencies above.

\paragraph{Cherry-count log-likelihood.}
The composite log-likelihood that the Maraschino fitter maximises is
\begin{align}
\mathcal{L}_{\text{cherry}}(\theta)
&= \mathcal{L}_{\text{singlet}}(\theta) + \sum_{b=1}^{n_\tau}
   \mathcal{L}_{\text{pair},b}(\theta), \\
\mathcal{L}_{\text{singlet}}(\theta)
&= \sum_{X,Y \in \alphabet \cup \{\sta,\fin\}}
   \Big(\textstyle\sum_b B^{(b)}_{XY}\Big)\,
   \log P^{\text{singlet}}_{XY}(\theta), \\
\mathcal{L}_{\text{pair},b}(\theta)
&= \sum_{u,v}\,
   \sum_{\anctok,\destok,\anctok',\destok'}
   C^{uv,(b)}_{\anctok\destok\,\anctok'\destok'}\,
   \log\!\bigg(\frac{F^{uv}_b(\anctok,\destok;\anctok',\destok')}
                       {Z^{u}_b(\anctok,\destok)}\bigg).
\label{eq:maraschino-pair-ll}
\end{align}
Here $P^{\text{singlet}}$ is the order-1 transition matrix obtained
from the MixDom Singlet HMM by row-normalising its adjacency frequencies
(Section~\ref{sec:order1hmm}); the singlet term scores the bigram
counts $B$ summed over time bins, since the singlet model is
time-independent. The pair term scores each per-bin adjacency tensor
under its own row-normalised conditional distribution.

\paragraph{Optimisation.}
$\mathcal{L}_{\text{cherry}}(\theta)$ is differentiable in
$\theta$. The maximisation is carried out by
gradient methods (Adam, optionally L-BFGS for refinement) on the
unconstrained parameterisation, using the same MixDom initialiser as
the exact Baum--Welch trainer (including the same flags for the number
of site classes, the class-equilibrium initialisation, and the
fragment-class assignment), so the two fitters can be started from
identical parameters and compared directly. Because the M-step of the
exact-EM trainer is closed-form and Maraschino's gradient optimiser
is not, this provides a controlled comparison of cherry-count fitting
against full Baum--Welch on the same data and architecture.

The fitted MixDom parameters are written to a checkpoint with the
same key layout as a \texttt{train\_pfam}-produced checkpoint, so that
either trainer's output can be loaded and refined by the other and
either can be used as input to the order-1 distillation
(Sections~\ref{sec:order1hmm} and \ref{sec:order1trans}).

\subsection{Distillation From MixDom To Order-1 Machines}
\label{sec:distillation-levels}

The MixDom HMMs defined above have large structured state spaces
(domains $\times$ fragments).
We now show how to distill these into compact \emph{order-1} machines---an HMM and a transducer---whose
transition probabilities depend only on the most recently emitted characters.
These machines are approximations of the full MixDom model.
The MixDom model parameters are assumed to have been trained, e.g.\ via Baum-Welch or via Maraschino cherry-count fitting (Section~\ref{sec:maraschino-fit}); see also \cite{LargeHolmes2026}.

The distillation operates on the alphabet $|\alphabet|$
(with the per-(domain, fragment-type) site-class mixture marginalised
implicitly inside the match emission tensors of
equation~(\ref{eq:maraschino-match-emprob})).
The Woodbury structural weights depend only on indel parameters;
the per-(class) emission tensors
$\eqm^{(\class)}_\anctok \exp(\revsub^{(\class)} \evoltime)_{\anctok\destok}$,
weighted by $\classdist_{\dom\frag\class}$, are contracted with those
weights to produce the order-1 transition probabilities.

Let $|\alphabet|$ denote the alphabet size and $\circ$ denote a distinguished
beginning-of-sequence (BOS) symbol.

\subsection{Notation for path marginalizations}

Both distillations require marginalizing over null (non-emitting) states
between consecutive emissions.
We introduce a compact notation for the expected frequency of partially observed paths.

In any HMM or transducer with states partitioned into emitting states $\emit$
and non-emitting (null) states $\silent$,
define the \emph{null closure}
\[
\nullcl \equiv (I - T_{\silent\silent})^{-1}
\]
where $T_{\silent\silent}$ is the submatrix of transitions among null states.
The effective transition matrix from emitting-or-start to emitting-or-end,
marginalizing all intervening null paths, is
\[
\effT_{ij} = T_{ij} + \sum_{p,q \in \silent} T_{ip}\, \nullcl_{pq}\, T_{qj}
\qquad i \in \emit \cup \{\sta\},\  j \in \emit \cup \{\fin\}
\]

For paths through the MixDom HMMs, we use the notation
\[
\pathE[\,\sta \to \underbrace{s_1}_{(\anctok_1,\destok_1)} \xrightarrow{\silent^\ast} \underbrace{s_2}_{(\anctok_2,\destok_2)} \to \cdots \to \fin\,]
\]
to denote the expected number of times the path visits the indicated sequence of emitting states
(with indicated emissions in subscript)
separated by zero or more null states (denoted $\silent^\ast$),
summed over all completions of the path to the left ($\sta \to \cdots$) and right ($\cdots \to \fin$).
Formally, if $\pi$ is the stationary distribution over states,
\[
\pathE[\, \underbrace{s_1}_{(\anctok_1,\destok_1)} \xrightarrow{\silent^\ast} \underbrace{s_2}_{(\anctok_2,\destok_2)} \,]
= \sum_{s_1 \in \emit} \sum_{s_2 \in \emit}
\left( \sum_{i} \pi_i \effT^\ast_{is_1} \right)
\emprob_{s_1}(\anctok_1,\destok_1)\,
\effT_{s_1 s_2}\,
\emprob_{s_2}(\anctok_2,\destok_2)
\left( \sum_{j} \effT^\ast_{s_2 j} \right)
\]
where $\emprob_s(\cdot)$ is the emission probability at state $s$
and $\effT^\ast_{ij} = [(I - \effT_{\emit\emit})^{-1}]_{ij}$ (with $\effT_{\emit\emit}$ the submatrix of $\effT$ restricted to emitting states) sums over all paths through emitting states.


\subsection{Distillation to Order-1 HMM}
\label{sec:distill-hmm}

\label{sec:order1hmm}

The MixDom Singlet HMM generates sequences from the stationary distribution.
We distill it into an order-1 HMM with states
\[
\{\sta, \fin\} \cup \{ \anctok : \anctok \in \alphabet \}
\]
where state $\anctok$ deterministically emits character $\anctok$.
Define the adjacency frequency from the full MixDom Singlet HMM:
\[
f(\anctok,\destok) =
\pathE[\,\underbrace{s_1}_{\anctok} \xrightarrow{\silent^\ast} \underbrace{s_2}_{\destok}\,]
\]
where the sum is over all emitting states $s_1,s_2$ weighted by emission probabilities
$\emprob_{s_1}(\anctok)$ and $\emprob_{s_2}(\destok)$ as above. Each
singlet emission probability $\emprob_{\ins_{\dom\frag}}(\anctok)
= \sum_\class \classdist_{\dom\frag\class} \eqm^{(\class)}_\anctok$
marginalises the per-(domain, fragment-type) class mixture.

\paragraph{Parameterization.}
The order-1 Singlet HMM transition probabilities are:
\begin{align}
P(\destok | \sta) &= \frac{\displaystyle\sum_{\destok'} f(\destok,\destok')}{\displaystyle\sum_{\anctok',\destok'} f(\anctok',\destok')}
&& \mbox{(start $\to$ first emission)} \label{eq:hmm-start} \\
P(\destok | \anctok) &= \frac{f(\anctok,\destok)}{\displaystyle\sum_{\destok'} f(\anctok,\destok')}
&& \mbox{(emission $\to$ emission)} \label{eq:hmm-trans} \\
P(\fin | \anctok) &= 1 - \sum_{\destok} P(\destok | \anctok)
&& \mbox{(emission $\to$ end)} \label{eq:hmm-end}
\end{align}
That is: normalize each row of the adjacency matrix $f$ to obtain transition probabilities,
allocating the residual probability mass to the end state.
The start distribution~(\ref{eq:hmm-start}) is proportional to the column marginals of $f$.


\subsection{Distillation to Order-1 WFST}

\label{sec:order1trans}

The MixDom Pair HMM describes the joint distribution over ancestor-descendant sequence pairs.
We distill it into an order-1 transducer (a Mealy machine)
whose state depends on the last inputted ancestor character
and the last outputted descendant character.

\paragraph{Machine states.}
The transducer has seven machine states $\{\sta,\mat,\ins,\del,\waitm,\waitd,\fin\}$
organized as a \emph{waiting machine}:
\begin{itemize}
\item \textbf{Non-waiting} (all outgoing transitions have $\varepsilon$ input):
  $\sta$ (start), $\mat$ (just matched), $\ins$ (just inserted), $\del$ (just deleted)
\item \textbf{Waiting} (all outgoing transitions consume input):
  $\waitm$ (ready after $\mat$ or $\ins$), $\waitd$ (ready after $\del$)
\item \textbf{Terminal}: $\fin$ (end)
\end{itemize}

The distinction $\waitm \neq \waitd$ is needed because outgoing transition weights differ:
in the MixDom Pair HMM, the $\mat/\ins$ rows use $\beta$ while $\del$ rows use $\gamma$.

\paragraph{Transitions.}
For a state with last-input ancestor $X$ and last-output descendant $Y$
(where $X,Y \in \alphabet \cup \{\circ\}$ and $\circ$ denotes BOS):

\medskip
\begin{tabular}{lllll}
Source & Dest & Input & Output & Weight \\
\hline
$\sta$ & $\waitm$ & $\varepsilon$ & $\varepsilon$ & $p_{\sta\waitm}$ \\
$\sta$ & $\ins$ & $\varepsilon$ & $\destok$ & $p_{\sta\ins}(Y,\destok)$ \\
$\sta$ & $\fin$ & $\varepsilon$ & $\varepsilon$ & $p_{\sta\fin}$ \\
\hline
$\waitm$ & $\mat$ & $\anctok$ & $\destok$ & $p_{\waitm\mat}(X,Y,\anctok,\destok)$ \\
$\waitm$ & $\del$ & $\anctok$ & $\varepsilon$ & $p_{\waitm\del}(X,Y,\anctok)$ \\
\hline
$\waitd$ & $\mat$ & $\anctok$ & $\destok$ & $p_{\waitd\mat}(X,Y,\anctok,\destok)$ \\
$\waitd$ & $\del$ & $\anctok$ & $\varepsilon$ & $p_{\waitd\del}(X,Y,\anctok)$ \\
\hline
$\mat$ & $\waitm$ & $\varepsilon$ & $\varepsilon$ & $p_{\mat\waitm}$ \\
$\mat$ & $\ins$ & $\varepsilon$ & $\destok$ & $p_{\mat\ins}(X,Y,\destok)$ \\
$\mat$ & $\fin$ & $\varepsilon$ & $\varepsilon$ & $p_{\mat\fin}$ \\
\hline
$\ins$ & $\waitm$ & $\varepsilon$ & $\varepsilon$ & $p_{\ins\waitm}$ \\
$\ins$ & $\ins$ & $\varepsilon$ & $\destok$ & $p_{\ins\ins}(X,Y,\destok)$ \\
$\ins$ & $\fin$ & $\varepsilon$ & $\varepsilon$ & $p_{\ins\fin}$ \\
\hline
$\del$ & $\waitd$ & $\varepsilon$ & $\varepsilon$ & $p_{\del\waitd}$ \\
$\del$ & $\ins$ & $\varepsilon$ & $\destok$ & $p_{\del\ins}(X,Y,\destok)$ \\
$\del$ & $\fin$ & $\varepsilon$ & $\varepsilon$ & $p_{\del\fin}$
\end{tabular}

\medskip
Note that $\waitm$ and $\waitd$ are \emph{not} associated with emissions;
they serve only to enforce the waiting-machine property.
Transitions from $\waitm$ and $\waitd$ always consume an ancestor input symbol;
transitions from $\mat$, $\ins$, $\del$, and $\sta$ never do.

\paragraph{Parameterization from the MixDom Pair HMM.}

We need the expected frequency of transitions conditioned on
(last ancestor input $X$, last descendant output $Y$, transition type, new symbols).
The key subtlety: through insert states, the last ancestor symbol $X$ must be \emph{propagated}
from the preceding match or delete, since inserts do not consume input.

Using the path notation from above, we enumerate all adjacency types
that arise in the MixDom Pair HMM.
Write $\mat[X,Y]$ for a match state that inputs $X$ and outputs $Y$,
$\ins[Y]$ for an insert that outputs $Y$,
and $\del[X]$ for a delete that inputs $X$.
The last-ancestor and last-descendant context is carried implicitly.

\medskip
\noindent
\textbf{Adjacency frequencies.}
The following table lists all adjacency types and their corresponding path marginalizations.
In each case, the frequency is computed as a sum over MixDom Pair HMM states,
with null states marginalized via $\nullcl$.
Write $\circ$ for the boundary (start/end) context.

\medskip
\begin{tabular}{llll}
Context & Adjacency & MixDom path & Frequency \\
\hline
\multicolumn{4}{l}{\emph{Start $\to$ Match:}} \\
& $\sta \to \mat[X',Y']$ &
$\sta \xrightarrow{\silent^\ast} \mat[X',Y']$ &
$f^{\sta\mat}(X',Y')$ \\
\multicolumn{4}{l}{\emph{Start $\to$ Insert:}} \\
& $\sta \to \ins[Y']$ &
$\sta \xrightarrow{\silent^\ast} \ins[Y']$ &
$f^{\sta\ins}(Y')$ \\
\multicolumn{4}{l}{\emph{Start $\to$ End (empty sequence):}} \\
& $\sta \to \fin$ &
$\sta \xrightarrow{\silent^\ast} \fin$ &
$f^{\sta\fin}$ \\
\hline
\multicolumn{4}{l}{\emph{Match $\to$ Match (via null states only):}} \\
& $\mat[X,Y] \to \mat[X',Y']$ &
$\mat[X,Y] \xrightarrow{\silent^\ast} \mat[X',Y']$ &
$f^{\mat\mat}(X,Y,X',Y')$ \\
\multicolumn{4}{l}{\emph{Match $\to$ Insert:}} \\
& $\mat[X,Y] \to \ins[Y']$ &
$\mat[X,Y] \xrightarrow{\silent^\ast} \ins[Y']$ &
$f^{\mat\ins}(X,Y,Y')$ \\
\multicolumn{4}{l}{\emph{Match $\to$ Delete:}} \\
& $\mat[X,Y] \to \del[X']$ &
$\mat[X,Y] \xrightarrow{\silent^\ast} \del[X']$ &
$f^{\mat\del}(X,Y,X')$ \\
\multicolumn{4}{l}{\emph{Match $\to$ End:}} \\
& $\mat[X,Y] \to \fin$ &
$\mat[X,Y] \xrightarrow{\silent^\ast} \fin$ &
$f^{\mat\fin}(X,Y)$ \\
\hline
\multicolumn{4}{l}{\emph{Insert $\to$ Insert (ancestor context $X$ propagated):}} \\
& $\ins[Y] \to \ins[Y']$ &
$\ins[Y] \xrightarrow{\silent^\ast} \ins[Y']$ &
$f^{\ins\ins}(X,Y,Y')$ \\
\multicolumn{4}{l}{\emph{Insert $\to$ Match (ancestor context $X$ propagated):}} \\
& $\ins[Y] \to \mat[X',Y']$ &
$\ins[Y] \xrightarrow{\silent^\ast} \mat[X',Y']$ &
$f^{\ins\mat}(X,Y,X',Y')$ \\
\multicolumn{4}{l}{\emph{Insert $\to$ Delete (ancestor context $X$ propagated):}} \\
& $\ins[Y] \to \del[X']$ &
$\ins[Y] \xrightarrow{\silent^\ast} \del[X']$ &
$f^{\ins\del}(X,Y,X')$ \\
\multicolumn{4}{l}{\emph{Insert $\to$ End (ancestor context $X$ propagated):}} \\
& $\ins[Y] \to \fin$ &
$\ins[Y] \xrightarrow{\silent^\ast} \fin$ &
$f^{\ins\fin}(X,Y)$ \\
\hline
\multicolumn{4}{l}{\emph{Delete $\to$ Match (descendant context $Y$ propagated):}} \\
& $\del[X] \to \mat[X',Y']$ &
$\del[X] \xrightarrow{\silent^\ast} \mat[X',Y']$ &
$f^{\del\mat}(X,Y,X',Y')$ \\
\multicolumn{4}{l}{\emph{Delete $\to$ Delete (descendant context $Y$ propagated):}} \\
& $\del[X] \to \del[X']$ &
$\del[X] \xrightarrow{\silent^\ast} \del[X']$ &
$f^{\del\del}(X,Y,X')$ \\
\multicolumn{4}{l}{\emph{Delete $\to$ Insert (descendant context $Y$ propagated):}} \\
& $\del[X] \to \ins[Y']$ &
$\del[X] \xrightarrow{\silent^\ast} \ins[Y']$ &
$f^{\del\ins}(X,Y,Y')$ \\
\multicolumn{4}{l}{\emph{Delete $\to$ End (descendant context $Y$ propagated):}} \\
& $\del[X] \to \fin$ &
$\del[X] \xrightarrow{\silent^\ast} \fin$ &
$f^{\del\fin}(X,Y)$
\end{tabular}

\medskip
In all cases, the ancestor context $X$ is propagated through insert states
(which do not consume input), and the descendant context $Y$ is propagated through delete states
(which do not emit output).
The notation $\xrightarrow{\silent^\ast}$ denotes zero or more transitions through null states,
marginalized via the null closure $\nullcl$.

\paragraph{Computing the frequencies.}
Each frequency above is computed from the MixDom Pair HMM as follows.
Let $\pi_s$ denote the stationary probability of state $s$
and $\effT_{s_1 s_2}$ the null-marginalized effective transition.
For the simplest case (direct adjacency):
\[
f^{\mat\del}(X,Y,X') = \sum_{s_1 \in \emit^{\mat}} \sum_{s_2 \in \emit^{\del}}
\left(\sum_i \pi_i \effT^\ast_{i s_1}\right)
\emprob_{s_1}(X,Y)\, \effT_{s_1 s_2}\, \emprob_{s_2}(X')
\left(\sum_j \effT^\ast_{s_2 j}\right)
\]
Boundary frequencies use the start/end rows of $\effT$:
\[
f^{\sta\mat}(X',Y') = \sum_{s \in \emit^{\mat}}
\effT_{\sta s}\, \emprob_s(X',Y')
\left(\sum_j \effT^\ast_{s j}\right),
\quad
f^{\mat\fin}(X,Y) = \sum_{s \in \emit^{\mat}}
\left(\sum_i \pi_i \effT^\ast_{i s}\right)
\emprob_s(X,Y)\, \effT_{s\,\fin}
\]
(and analogously for $f^{\sta\ins}$, $f^{\sta\fin}$, $f^{\ins\fin}$, $f^{\del\fin}$).

\paragraph{Normalization to transducer parameters.}
Given the adjacency frequencies, the order-1 transducer weights are obtained by normalization.
For each context $(X,Y)$ and source machine state, normalize outgoing weights to sum to 1:

\medskip
\noindent\emph{Start transitions} (context $\circ$):
\begin{align*}
p_{\sta\waitm} &= \frac{\sum_{X',Y'} f^{\sta\mat}(X',Y')}
  {\sum_{X',Y'} f^{\sta\mat}(X',Y') + \sum_{Y'} f^{\sta\ins}(Y') + f^{\sta\fin}}, \quad
p_{\sta\ins}(\destok) = \frac{f^{\sta\ins}(\destok)}
  {\sum_{X',Y'} f^{\sta\mat}(X',Y') + \sum_{Y'} f^{\sta\ins}(Y') + f^{\sta\fin}}
\end{align*}
(and $p_{\sta\fin}$ uses the same denominator with numerator $f^{\sta\fin}$).

\medskip
\noindent\emph{Wait-after-match/insert transitions} (context $(X,Y)$, consuming input $\anctok$):
\begin{align*}
p_{\waitm\mat}(X,Y,\anctok,\destok) &= \frac{f^{\cdot\mat}(X,Y,\anctok,\destok)}
  {\sum_{\anctok'}\left[\sum_{\destok'} f^{\cdot\mat}(X,Y,\anctok',\destok') + f^{\cdot\del}(X,Y,\anctok')\right]} \\
p_{\waitm\del}(X,Y,\anctok) &= \frac{f^{\cdot\del}(X,Y,\anctok)}
  {\sum_{\anctok'}\left[\sum_{\destok'} f^{\cdot\mat}(X,Y,\anctok',\destok') + f^{\cdot\del}(X,Y,\anctok')\right]}
\end{align*}
where $f^{\cdot\mat}(X,Y,\anctok,\destok) = f^{\mat\mat}(X,Y,\anctok,\destok) + f^{\ins\mat}(X,Y,\anctok,\destok)$
and $f^{\cdot\del}(X,Y,\anctok) = f^{\mat\del}(X,Y,\anctok) + f^{\ins\del}(X,Y,\anctok)$,
aggregating over both match and insert sources that share the $\waitm$ wait state.

\medskip
\noindent\emph{Wait-after-delete transitions} (context $(X,Y)$, consuming input $\anctok$):
\begin{align*}
p_{\waitd\mat}(X,Y,\anctok,\destok) &= \frac{f^{\del\mat}(X,Y,\anctok,\destok)}
  {\sum_{\anctok'}\left[\sum_{\destok'} f^{\del\mat}(X,Y,\anctok',\destok') + f^{\del\del}(X,Y,\anctok')\right]} \\
p_{\waitd\del}(X,Y,\anctok) &= \frac{f^{\del\del}(X,Y,\anctok)}
  {\sum_{\anctok'}\left[\sum_{\destok'} f^{\del\mat}(X,Y,\anctok',\destok') + f^{\del\del}(X,Y,\anctok')\right]}
\end{align*}

\medskip
\noindent\emph{Post-match transitions} (context $(X,Y)$, after matching with $\anctok,\destok$;
new context becomes $(\anctok,\destok)$):
\begin{align*}
p_{\mat\waitm}(X,Y) &= \frac{\sum_{\anctok',\destok'} f^{\mat\mat}(X,Y,\anctok',\destok') + \sum_{\anctok'} f^{\mat\del}(X,Y,\anctok')}
  {Z^{\mat}(X,Y)} \\
p_{\mat\ins}(X,Y,\destok) &= \frac{f^{\mat\ins}(X,Y,\destok)}{Z^{\mat}(X,Y)}, \quad
p_{\mat\fin}(X,Y) = \frac{f^{\mat\fin}(X,Y)}{Z^{\mat}(X,Y)}
\end{align*}
where $Z^{\mat}(X,Y) = \sum_{\anctok',\destok'} f^{\mat\mat}(X,Y,\anctok',\destok') + \sum_{\anctok'} f^{\mat\del}(X,Y,\anctok') + \sum_{\destok} f^{\mat\ins}(X,Y,\destok) + f^{\mat\fin}(X,Y)$.

\medskip
\noindent\emph{Post-insert transitions} are analogous, using $f^{\ins\cdot}$ frequencies with $Z^{\ins}(X,Y)$.

\medskip
\noindent\emph{Post-delete transitions} (context $(X,Y)$, after deleting $\anctok$;
new context becomes $(\anctok,Y)$):
\begin{align*}
p_{\del\waitd}(X,Y) &= \frac{\sum_{\anctok'} f^{\del\del}(X,Y,\anctok') + \sum_{\anctok',\destok'} f^{\del\mat}(X,Y,\anctok',\destok')}
  {Z^{\del}(X,Y)} \\
p_{\del\ins}(X,Y,\destok) &= \frac{f^{\del\ins}(X,Y,\destok)}{Z^{\del}(X,Y)}, \quad
p_{\del\fin}(X,Y) = \frac{f^{\del\fin}(X,Y)}{Z^{\del}(X,Y)}
\end{align*}
where $Z^{\del}(X,Y) = \sum_{\anctok'} f^{\del\del}(X,Y,\anctok') + \sum_{\anctok',\destok'} f^{\del\mat}(X,Y,\anctok',\destok') + \sum_{\destok} f^{\del\ins}(X,Y,\destok) + f^{\del\fin}(X,Y)$.
